\documentclass[11pt,a4paper]{article}
\usepackage[T1]{fontenc}
\usepackage[utf8]{inputenc}
\usepackage{lmodern}
\usepackage[margin=2.3cm]{geometry}
\usepackage{booktabs}
\usepackage{tabularx}
\usepackage{longtable}
\usepackage{array}
\usepackage[dvipsnames]{xcolor}
\usepackage{enumitem}
\usepackage{titlesec}
\usepackage{parskip}
\usepackage{setspace}
\usepackage{needspace}
\usepackage[most]{tcolorbox}
\usepackage[hidelinks]{hyperref}

\setstretch{1.25}
\setlength{\parskip}{0.8em}
\setlist{itemsep=0.35em, topsep=0.4em}
\renewcommand{\arraystretch}{1.45}
\definecolor{accent}{RGB}{30,80,150}
\titleformat{\section}{\Large\bfseries\color{accent}}{\thesection}{0.8em}{}
\titleformat{\subsection}{\large\bfseries}{\thesubsection}{0.8em}{}
\titlespacing*{\section}{0pt}{2.0em}{1.0em}
\titlespacing*{\subsection}{0pt}{1.4em}{0.6em}
\newcolumntype{L}{>{\raggedright\arraybackslash}X}
\newcolumntype{P}[1]{>{\raggedright\arraybackslash}p{#1}}
\newcommand{\file}[1]{\texttt{\small #1}}
\newcommand{\keybox}[1]{\par\smallskip\noindent\colorbox{accent!8}{\parbox{\dimexpr\linewidth-2\fboxsep}{\vspace{0.5em}#1\vspace{0.5em}}}\par\medskip}

% ---- family colours (the same as the dashboard)
\definecolor{cFlood}{HTML}{EF4444}
\definecolor{cDDoS}{HTML}{BE123C}
\definecolor{cSpoofFlood}{HTML}{8B5CF6}
\definecolor{cFrag}{HTML}{EA580C}
\definecolor{cRecon}{HTML}{0891B2}
\definecolor{cBrute}{HTML}{D97706}
\definecolor{cARP}{HTML}{7C3AED}
\definecolor{cNone}{HTML}{64748B}
% ---- badge colours
\definecolor{bModel}{HTML}{2563EB}
\definecolor{bRule}{HTML}{16A34A}
\definecolor{bKernel}{HTML}{475569}
\definecolor{bLive}{HTML}{15803D}
\definecolor{bDocker}{HTML}{0E7490}
\definecolor{bData}{HTML}{B45309}
\definecolor{bNo}{HTML}{B91C1C}

\newcommand{\badge}[2]{\tcbox[on line, colback=#1, colframe=#1, arc=2pt, boxsep=1.5pt, left=3pt, right=3pt, top=1pt,
  bottom=1pt]{\textcolor{white}{\scriptsize\bfseries #2}}}
\newcommand{\MODEL}{\badge{bModel}{MODEL}}
\newcommand{\RULE}[1]{\badge{bRule}{RULE: #1}}
\newcommand{\KERNEL}{\badge{bKernel}{LINUX KERNEL}}
\newcommand{\LIVE}{\badge{bLive}{TESTED LIVE ON THE PI}}
\newcommand{\DOCKER}{\badge{bDocker}{TESTED IN DOCKER}}
\newcommand{\DATA}{\badge{bData}{DATASET ROWS ONLY}}
\newcommand{\NOTEST}{\badge{bNo}{NOT TESTED}}
\newcommand{\NOTCOV}{\badge{bNo}{NOT COVERED}}

\newtcolorbox{attack}[2]{enhanced, breakable, colback=#1!3, colframe=#1!80!black, coltitle=white,
  colbacktitle=#1!85!black, title={\large\bfseries #2}, boxrule=0.8pt, arc=4pt, left=8pt, right=8pt, top=6pt,
  bottom=6pt, before skip=1.4em, after skip=1.4em}
\newcommand{\fld}[2]{\par\vspace{0.55em}\noindent{\bfseries\color{black!75}#1}\par\vspace{-0.4em}\noindent #2\par}

\begin{document}

\begin{titlepage}
\centering
\vspace*{4.5cm}
{\Huge\bfseries\color{accent} TinyIDS\par}
\vspace{1.2em}
{\LARGE The attacks: what they are and how we stop them\par}
\vspace{2.5em}
{\large Every attack the Raspberry Pi gateway detects and mitigates: what it does, why attackers use it, the harm it
causes, how it is usually stopped, and how TinyIDS stops it\par}
\vfill
{\large 18 November 2026\par}
\vspace{2cm}
\end{titlepage}

\tableofcontents
\newpage

% ===============================================================================================================
\section{Viva: problem, layers and defence}

\subsection{The problem statement}

\keybox{IoT devices (here: NodeMCU sensors) are cheap and have almost no security of their own: no firewall, weak
passwords, no updates. They all talk through one \textbf{gateway} (the Raspberry Pi). We build a small
\textbf{intrusion detection and prevention system} that runs \emph{on that gateway}, in real time, on a
Raspberry Pi 4: it watches every packet that reaches the gateway, detects floods (DoS / DDoS), scans, MQTT password
guessing and spoofing, \textbf{blocks} the attacker automatically, and \textbf{explains} every alert with its top-3
reasons (SHAP).}

\subsection{The layers in one table}

\begin{center}
\begin{tabularx}{\linewidth}{P{2.5cm}P{4.3cm}P{3.0cm}L}
\toprule
Layer & Its job & What it adds & Example in our system \\
\midrule
7 Application & what the programs say to each other & the message & MQTT, HTTP (dashboard) \\
6 Presentation & format, encryption & (TLS, JSON) & JSON sensor messages \\
5 Session & keep a conversation going & -- & the MQTT session \\
4 Transport & from program to program; reliable (TCP) or fast (UDP) & \textbf{port numbers}, TCP flags & TCP 1884,
UDP 5683, SYN/ACK/FIN \\
3 Network & from device to device, across networks & \textbf{IP addresses} & 172.29.3.125 (Pi) \\
2 Data link & to the next device on the same network & \textbf{MAC addresses} & Wi-Fi frames, ARP \\
1 Physical & bits as signals & -- & the Wi-Fi radio \\
\bottomrule
\end{tabularx}
\end{center}

Going down the stack (sending), each layer wraps the data with its own header: message $\rightarrow$
+ port (4) $\rightarrow$ + IP (3) $\rightarrow$ + MAC (2) $\rightarrow$ radio (1). The receiver unwraps it the other
way. \textbf{ARP} links layer 3 to 2: it answers ``which MAC has this IP?''. In the TCP/IP model, layers 5--7 are
one ``application'' layer.

\subsection{Which layer are we securing?}

\textbf{Mainly the network and transport layers (OSI layers 3 and 4)} of the traffic arriving at the gateway: the
model reads IP / TCP / UDP / ICMP header information (protocol, TCP flags, packet sizes). Around that, our rules also
cover \textbf{layer 2} (ARP spoofing: MAC $\leftrightarrow$ IP) and \textbf{layer 7} (MQTT brute force: the broker's
login answers).

In the 3-layer IoT model: the NodeMCUs are the \textbf{perception layer}, the Wi-Fi and the Pi gateway are the
\textbf{network layer} (that is where TinyIDS sits), and the broker and dashboard are the \textbf{application layer}.

\textbf{Type of IDS:} a \textbf{network-based IDS} placed on the gateway (it watches traffic, not files or processes),
that also acts as an \textbf{IPS} (it blocks, through the gateway's firewall). Detection = \textbf{machine learning}
(a supervised classifier trained on labelled attacks) + \textbf{rules} (signature-like thresholds).

\subsection{Every part of the system and its layer}

\begin{longtable}{P{3.5cm}P{3.6cm}P{7.9cm}}
\toprule
Part & Layer (OSI) & Why \\
\midrule
\endhead
Wi-Fi (IIITA) & 1 physical, 2 data link & radio + 802.11 frames with MAC addresses; everything travels over it \\
NodeMCU (ESP8266) & all layers 1--7 (an end device) & Wi-Fi radio (1--2), gets an IP (3), opens TCP to the broker (4),
speaks MQTT (7). In IoT terms: perception layer \\
MQTT & \textbf{7 application}, over TCP (4) & publish / subscribe protocol, TCP port 1884 (ours) \\
Mosquitto broker (Pi) & 7 application & a server program that passes MQTT messages on \\
Raspberry Pi (gateway) & 3 network device; runs all layers & every device sends to it; it captures, decides,
blocks \\
Packet capture (scapy) & reads at 2 (frames) & sees the whole Ethernet / Wi-Fi frame incl. MAC addresses \\
\textbf{The ML model} & \textbf{uses layers 3--4} information & features = IP protocol, TCP flags, packet sizes;
the program itself runs as an application on the Pi \\
PortScanRule & 4 transport & counts different destination \textbf{ports} \\
FloodCounter & 3 network & counts different source \textbf{IP addresses} (DoS vs DDoS) \\
BruteForceRule & 7 application & reads the broker's MQTT CONNACK answers (refused logins) \\
AddressWatcher & 2 data link / 3 network & checks MAC $\leftrightarrow$ IP pairs (ARP and IP spoofing) \\
\textbf{Mitigation (iptables)} & \textbf{3--4} (and 2 for MAC blocks) & Linux kernel firewall (netfilter): drops by
source IP, rate-limits packets, drops by MAC; SYN cookies work at 4 \\
Dashboard (Streamlit) & 7 application & a web page (HTTP) that is also an MQTT client \\
Laptop & an end device (all layers) & operator (dashboard, SSH) and attacker \\
WSL & not a network layer & just Linux running on the laptop, sharing its Wi-Fi card, so Linux attack tools run
there \\
hping3 & builds packets at 3--4 & hand-made IP / TCP / UDP / ICMP packets (floods) \\
nmap & 4 (ports), 3 (OS guess), 7 (versions) & port scan, OS fingerprint, service versions \\
mosquitto\_pub (brute force) & 7 application & real MQTT logins with wrong passwords \\
arpspoof & 2 data link & fake ARP replies \\
bot.py (bot master) & 4 (UDP), message at 7 & one UDP packet with the text ``udp 10'' to port 6667 \\
\bottomrule
\end{longtable}

\subsection{The attacks by layer}

\begin{center}
\begin{tabularx}{\linewidth}{P{2.6cm}P{5.0cm}L}
\toprule
Layer & Attacks & Our defence \\
\midrule
2 data link & ARP spoofing & AddressWatcher: block the MAC, fix the ARP table (on the college Wi-Fi the access point
already rewrites spoofed ARP) \\
3 network & ICMP flood, IP-spoofed floods, fragmentation floods, DDoS (many sources) & model (Flood) + FloodCounter
(DoS vs DDoS) + AddressWatcher (spoofed sources) $\rightarrow$ iptables block \\
4 transport & SYN, ACK, PSH-ACK, RST-FIN and UDP floods; port scan & model (Flood / Recon) + PortScanRule + SYN
cookies $\rightarrow$ iptables block / rate-limit \\
7 application & MQTT brute force (version / OS scans touch 7 and 3 too) & BruteForceRule $\rightarrow$ block;
web attacks, HTTP floods, DNS spoofing: out of scope \\
\bottomrule
\end{tabularx}
\end{center}

\subsection{How we mitigate, in one paragraph}

Detection at the gateway, \textbf{enforcement in the gateway's Linux kernel}: TinyIDS adds an \textbf{iptables}
rule, so the kernel drops the attacker's packets before they reach the broker. Sure attacks are \textbf{blocked}
(30 s), our own compromised NodeMCUs are \textbf{quarantined} (120 s), unsure ones are \textbf{rate-limited}
(20 packets/s), spoofers are blocked \textbf{by MAC}. Every action expires by itself, the operator's laptop is never
locked out, and the dashboard can switch blocking off or release everything.

\subsection{Likely questions, short answers}

\begin{longtable}{P{5.2cm}P{9.8cm}}
\toprule
Question & Answer \\
\midrule
\endhead
Why at the gateway, not on each sensor? & A NodeMCU has about 80 KB of RAM and no firewall; the gateway sees all traffic
and can block it in one place. \\
Why not just use TLS / passwords? & TLS protects \emph{confidentiality}: nobody can read the messages. A flood
attacks \emph{availability}: encryption does not stop 1000 packets a second. Passwords are there (broker login), and
the brute-force rule protects them. \\
DoS vs DDoS? & DoS = one attacker; DDoS = many at once (a botnet). The model says ``Flood'', the FloodCounter rule
counts the senders (2+ in 10 s = DDoS). \\
What if the attacker fakes his IP? & Then blocking the IP is useless: the AddressWatcher sees one MAC using many IPs
and blocks the MAC. \\
Signature-based or anomaly-based? & Neither purely: a \emph{supervised} ML classifier trained on labelled attacks,
plus threshold rules. Not anomaly detection (it does not learn ``normal'' only). \\
Inline or passive? & It sniffs a copy of every packet (passive detection) and enforces through the kernel firewall
(active prevention): an IDS + IPS. \\
False positives? & Threshold 0.70, 2 windows in a row for unsure cases, the admin laptop never blocked, a 30 s grace
after a quarantine. Live: 0 false alarms on the 3 NodeMCUs. \\
What can it not stop? & Web attacks (no web server), low-rate DDoS, attacks hidden inside encrypted traffic, and
anything at layer 1 (radio jamming). \\
\bottomrule
\end{longtable}

\newpage
\section{Demo cheat sheet: command $\rightarrow$ what the dashboard shows}

For every attack: the command that starts it, where to type it, the alert name the dashboard shows, and the
\textbf{top-3 SHAP reasons} (measured live on the Pi, 17 November). \file{PI} = the Pi's IP address.
All floods run for a fixed time: put \texttt{timeout -s INT 15} in front (stops after 15 s). The scripts do this
for you: \texttt{bash attack\_all.sh PI} (every attack in a row) or \texttt{python3 launch.py PI} (a menu).

\newcommand{\cmd}[1]{\texttt{\scriptsize #1}}
\newcommand{\where}[1]{\par{\scriptsize\color{black!60}#1}}
\newcommand{\shap}[1]{\textcolor{bModel}{\texttt{\scriptsize #1}}}
\newcommand{\rul}[1]{\textcolor{bRule}{\texttt{\scriptsize #1}}}

\begin{longtable}{P{3.1cm}P{6.1cm}P{5.9cm}}
\toprule
Attack & Command & Dashboard alert \;/\; top-3 reasons \\
\midrule
\endhead
\multicolumn{3}{l}{\textcolor{cFlood}{\textbf{Floods (DoS), from the laptop}}} \\
SYN flood \where{WSL} & \cmd{sudo hping3 -S -p 1883 -i u1000 PI} & \textbf{DoS flood} 0.96 \newline
\shap{Min, Std, ack\_flag\_number} \newline {\scriptsize tiny equal packets, no ACK flag} \\
UDP flood \where{WSL} & \cmd{sudo hping3 --udp -p 5683 -d 8 -i u1000 PI} & \textbf{DoS flood} 1.00 \newline
\shap{Protocol Type, Min, UDP} \newline {\scriptsize protocol 17 = UDP} \\
ICMP flood \where{WSL} & \cmd{sudo hping3 --icmp -i u1000 PI} & \textbf{DoS flood} 1.00 \newline
\shap{ICMP, Min, ack\_flag\_number} \newline {\scriptsize all packets are pings} \\
ACK flood \where{WSL} & \cmd{sudo hping3 -A -p 1883 -i u1000 PI} & \textbf{DoS flood} 0.98 \newline
\shap{Min, Max, Std} \newline {\scriptsize every packet exactly the same tiny size} \\
PSH-ACK flood \where{WSL} & \cmd{sudo hping3 -PA -p 1883 -i u1000 PI} & \textbf{DoS flood} 1.00 \newline
\shap{psh\_flag\_number, Min, Max} \newline {\scriptsize PSH flag in every packet} \\
RST-FIN flood \where{WSL} & \cmd{sudo hping3 -RF -p 1883 -i u1000 PI} & \textbf{DoS flood} 1.00 \newline
\shap{fin\_flag\_number, Min, rst\_flag\_number} \newline {\scriptsize FIN + RST in every packet} \\
\midrule
\multicolumn{3}{l}{\textcolor{cDDoS}{\textbf{DDoS, from the 3 NodeMCU bots}}} \\
Bot TCP DDoS \where{Windows (PowerShell), in \file{laptop/}} & \cmd{python bot.py tcp 10} &
\textbf{DDoS flood (3 senders)} 0.96--0.97, all 3 \textbf{quarantined} \newline
\shap{rst\_flag\_number, syn\_flag\_number, Min} \newline {\scriptsize open + close connections} \\
Bot UDP DDoS \where{Windows (PowerShell)} & \cmd{python bot.py udp 10} & \textbf{DDoS flood (3 senders)} 1.00,
\textbf{quarantined} \newline \shap{Min, Protocol Type, UDP} \\
Stop the bots \where{Windows} & \cmd{python bot.py stop} & -- \\
\midrule
\multicolumn{3}{l}{\textcolor{cRecon}{\textbf{Scans, from the laptop}}} \\
Port scan \where{WSL} & \cmd{nmap -sT -T4 -p 1-1024 PI} & \textbf{Recon} 0.99 (model) \newline
\shap{AVG, Tot size, Max} \newline + \textbf{Recon} 1.00 (rule) \newline
\rul{different ports, within seconds, rule threshold} \\
Version scan \where{WSL} & \cmd{nmap -sV --top-ports 100 PI} & \textbf{Recon} 0.99 \newline
\shap{AVG, Tot size, Max} \\
OS scan \where{WSL} & \cmd{sudo nmap -O PI} & \textbf{DoS flood} 0.96 (wrong name, still blocked) \newline
\shap{Min, Std, syn\_flag\_number} \\
Vulnerability scan \where{WSL (menu only)} & \cmd{nmap --script vuln -p 1884 PI} & expected \textbf{Recon}
{\scriptsize (not in the final live run)} \\
\midrule
\multicolumn{3}{l}{\textcolor{cBrute}{\textbf{Password attack, from the laptop}}} \\
MQTT brute force \where{WSL} & \cmd{for p in 123456 password admin 12345678 qwerty; do timeout 3 mosquitto\_pub -h
PI -p 1884 -u tinyids -P \$p -t test -m hi; done} & \textbf{BruteForce} 1.00 (rule) \newline
\rul{refused MQTT logins, within seconds, rule threshold} \\
\midrule
\multicolumn{3}{l}{\textcolor{cSpoofFlood}{\textbf{Spoofing (tested in Docker only)}}} \\
IP-spoofed flood \where{WSL (from the old project's \file{launch.py})} &
\cmd{sudo hping3 -S -p 1884 -i u2000 --rand-source PI} \newline
\cmd{sudo hping3 --udp -p 5683 -d 32 -i u2000 --rand-source PI} & \textbf{IP spoofing} (rule), MAC blocked
\newline \rul{source IPs from one MAC, within seconds, rule threshold} \\
ARP spoofing \where{WSL, menu option 7} & \cmd{sudo arpspoof -i IFACE -t PI NODEMCU\_IP} & \textbf{ARP spoofing}
(rule), MAC blocked + ARP fixed \newline \rul{IP claimed, new MAC, known MAC} \\
\midrule
Normal traffic & 3 NodeMCUs on, no attack & \textbf{no alert} (0 false alarms live) \\
\bottomrule
\end{longtable}

\textbf{Model or rule?} Every alert says where it came from (\file{source}: \texttt{model} or \texttt{rule}).
\shap{Blue} reasons are SHAP features (the model); \rul{green} reasons are the counts of a rule (no SHAP, a rule
has no model to explain).

\subsection*{Reading the SHAP reasons in 5 seconds}

\begin{center}
\begin{tabularx}{\linewidth}{P{5.6cm}L}
\toprule
If the top reasons contain\ldots & it is most likely\ldots \\
\midrule
\shap{Min}, \shap{Std}, \shap{ack\_flag\_number} (tiny packets, no ACK) & SYN flood (live) \\
\shap{syn\_flag\_number} & SYN packets: bot TCP flood, or an OS scan (named ``Flood'' live) \\
\shap{UDP}, \shap{Protocol Type} & UDP flood (protocol 17) \\
\shap{ICMP} & ping flood \\
\shap{psh\_flag\_number} & PSH-ACK flood \\
\shap{rst\_flag\_number}, \shap{fin\_flag\_number} & RST-FIN flood, or bot TCP flood (connections closed) \\
only sizes: \shap{Min}, \shap{Max}, \shap{Std} & ACK flood (identical tiny packets) \\
\shap{AVG}, \shap{Tot size}, \shap{Max} with \textbf{Recon} & a scan (larger, varied packets) \\
``(N senders)'' in the alert name & DDoS: the FloodCounter rule counted N flooding senders \\
\bottomrule
\end{tabularx}
\end{center}

\newpage
\section{What every part of the commands means}

\subsection{First: the TCP flags, in plain words}

Every TCP packet carries small on/off switches called \textbf{flags}. They say what the packet is for. A normal
conversation uses them like a phone call:

\begin{center}
\begin{tabularx}{\linewidth}{P{2.0cm}P{4.4cm}L}
\toprule
Flag & Means & In a phone call \\
\midrule
\textbf{SYN} & ``I want to start a connection'' & dialling the number \\
\textbf{ACK} & ``I received what you sent'' & ``mm-hm, I heard you'' \\
\textbf{PSH} & ``here is data, hand it to the program now'' & ``listen, this is important'' \\
\textbf{FIN} & ``I am finished, let us close politely'' & ``bye!'' \\
\textbf{RST} & ``stop at once, this connection is cancelled'' & slamming the phone down \\
\bottomrule
\end{tabularx}
\end{center}

A normal connection: SYN $\rightarrow$ SYN+ACK $\rightarrow$ ACK, then data with PSH+ACK, then FIN to close. A flood
sends one of these over and over, thousands of times a second, with no real conversation behind it.

\subsection{hping3: the flood tool}

\textbf{hping3} builds packets exactly the way you describe them and fires them at a target. Example:
\cmd{timeout -s INT 15 sudo hping3 -S -p 1883 -i u1000 PI}

\begin{longtable}{P{3.6cm}P{11.4cm}}
\toprule
Part & What it means \\
\midrule
\endhead
\cmd{timeout -s INT 15} & ``stop the following command after 15 seconds''. \cmd{-s INT} = stop it the same way
as pressing Ctrl+C, so hping3 still prints its summary (how many packets it sent). \\
\cmd{sudo} & ``run as administrator''. Building your own raw packets needs admin rights. \\
\cmd{hping3} & the program itself. \\
\cmd{-S} & set the \textbf{SYN} flag in every packet $\rightarrow$ SYN flood. \\
\cmd{-A} & set the \textbf{ACK} flag $\rightarrow$ ACK flood. \\
\cmd{-PA} & two flags together: \cmd{-P} = \textbf{PSH} and \cmd{-A} = \textbf{ACK} $\rightarrow$ PSH-ACK flood. \\
\cmd{-RF} & two flags together: \cmd{-R} = \textbf{RST} and \cmd{-F} = \textbf{FIN} $\rightarrow$ RST-FIN flood. \\
\cmd{--udp} & send UDP packets instead of TCP (UDP has no flags and no connection). \\
\cmd{--icmp} & send ICMP packets: ping requests. \\
\cmd{-p 1883} & the \textbf{port}: which ``door'' of the Pi to hit. 1883 = the lab's MQTT broker, 1884 = ours,
5683 = the CoAP port used for the UDP floods. \\
\cmd{-i u1000} & the \textbf{interval} between two packets: \cmd{u} = microseconds, so 1000 microseconds = 1
millisecond $\rightarrow$ about \textbf{1000 packets every second}. (\cmd{u2000} = about 500 per second.) \\
\cmd{-d 8} & put \textbf{8 bytes} of data in every packet (default: none). \cmd{-d 32} = 32 bytes. \\
\cmd{--rand-source} & write a \textbf{random fake sender address} on every packet (IP spoofing). \\
\cmd{-c 3} & send only \textbf{3} packets, then stop (used by \file{attack\_all.sh} to check that the Pi can be
reached). \\
\cmd{-f} & cut every packet into fragments (a fragmentation flood; we never used it). \\
\cmd{PI} & the target: write the Pi's IP address here, e.g.\ \cmd{172.29.2.187}. \\
\bottomrule
\end{longtable}

\subsection{nmap: the scanning tool}

\textbf{nmap} knocks on the target's ports to find out which are open, and what is running behind them.

\begin{longtable}{P{3.6cm}P{11.4cm}}
\toprule
Part & What it means \\
\midrule
\endhead
\cmd{-sT} & \textbf{connect scan}: for every port, try a normal, full connection. If it connects, the port is open.
Needs no admin rights. \\
\cmd{-T4} & \textbf{speed} setting from 0 (very slow, sneaky) to 5 (very fast). 4 = fast. \\
\cmd{-p 1-1024} & which \textbf{ports} to try: number 1 up to 1024 (where the common services live). \\
\cmd{-sV} & \textbf{version scan}: on every open port, talk to the program a little until it tells its name and
version (e.g.\ ``Mosquitto 2.0''). \\
\cmd{--top-ports 100} & try only the 100 most commonly used ports (faster than all of them). \\
\cmd{-O} & \textbf{OS scan}: send unusual packets and guess the operating system from how it answers. Needs
\cmd{sudo}. \\
\cmd{--script vuln} & run nmap's ready-made scripts that test for \textbf{known weaknesses}
(vulnerabilities). \\
\bottomrule
\end{longtable}

\subsection{mosquitto\_pub: the brute-force attack}

\cmd{for p in 123456 password admin; do timeout 3 mosquitto\_pub -h PI -p 1884 -u tinyids -P \$p -t test -m hi; done}

\begin{longtable}{P{3.6cm}P{11.4cm}}
\toprule
Part & What it means \\
\midrule
\endhead
\cmd{for p in A B C; do \ldots; done} & a \textbf{loop}: do the same thing once with \cmd{p} = A, then with
\cmd{p} = B, then C. Here: once for every password in the list. \\
\cmd{mosquitto\_pub} & a small program that logs in to an MQTT broker and publishes one message. \\
\cmd{-h PI} & \textbf{host}: the broker's address (the Pi). \\
\cmd{-p 1884} & the broker's \textbf{port} (ours). \\
\cmd{-u tinyids} & the \textbf{user name} to log in with. \\
\cmd{-P \$p} & the \textbf{password} to try: \cmd{\$p} is replaced by the current password from the list. \\
\cmd{-t test} & the \textbf{topic} to publish to. \\
\cmd{-m hi} & the \textbf{message} to publish. (It never arrives: the password is wrong, the broker refuses.) \\
\cmd{timeout 3} & give up on one try after 3 seconds. \\
\bottomrule
\end{longtable}

\subsection{The other commands}

\begin{longtable}{P{4.6cm}P{10.4cm}}
\toprule
Command & What it means \\
\midrule
\endhead
\cmd{python bot.py tcp 10} & run the bot master: tell every NodeMCU bot to do a \textbf{TCP} flood for \textbf{10
seconds}. \cmd{udp 10} = UDP flood, \cmd{stop} = stop now. (Windows, in the \file{laptop} folder.) \\
\cmd{sudo arpspoof -i IFACE -t PI NODEMCU\_IP} & ARP spoofing. \cmd{-i IFACE} = which network card of the laptop
to send from; \cmd{-t PI} = the device to fool (the \textbf{target}); \cmd{NODEMCU\_IP} = the address to pretend
to be. It keeps telling the Pi ``this NodeMCU's address is at \emph{my} network card''. \\
\cmd{bash attack\_all.sh PI} & run every attack one after another, with pauses, and save the logs. \\
\cmd{python3 launch.py PI} & a menu: pick one attack by number. \\
\cmd{./run\_ids.sh admin LAPTOP\_IP} & on the Pi: start our broker and the detector. \cmd{admin LAPTOP\_IP} = the
operator's laptop: its SSH and dashboard connection are never blocked. \\
\cmd{\ldots\ | tee lab\_log.txt} & \cmd{|} (a ``pipe'') hands the output of the left command to the right one.
\cmd{tee} then shows it on the screen \emph{and} writes the same text into \file{lab\_log.txt} (like a T-shaped
pipe: one input, two outputs). \\
\cmd{mosquitto\_sub -t ids/alerts -v} & listen to a topic and print every message that arrives (the alert logger
in \file{attack\_all.sh}). \cmd{-v} = also print the topic name. \\
\bottomrule
\end{longtable}

\newpage
\section{How to read this document}

Every attack has its own card in its family's colour. Each card answers the same questions: what it is, how it
works, why an attacker does it, what harm it does, how it is usually stopped, and how TinyIDS detects and stops it.

\subsection{Who detects it}

\begin{center}
\begin{tabularx}{\linewidth}{P{3.6cm}L}
\toprule
Badge & Meaning \\
\midrule
\MODEL & the LightGBM model (Benign / Flood / Recon), on every window of 10 packets from one sender \\
\RULE{name} & a simple rule written by us (no machine learning); the name says which one \\
\KERNEL & Linux itself (SYN cookies, which the start script turns on) \\
\bottomrule
\end{tabularx}
\end{center}

\subsection{How well we know it works}

\begin{center}
\begin{tabularx}{\linewidth}{P{5.0cm}L}
\toprule
Badge & Meaning \\
\midrule
\LIVE & real attack against the real Pi in the lab (17 November), detected and blocked \\
\DOCKER & tested in our Docker copy of the testbed (same detector code), not yet on the real Pi \\
\DATA & never sent on our testbed; the number is how often the model calls the dataset's test rows of it an attack \\
\NOTEST & the protection exists, but this attack was never tried \\
\NOTCOV & TinyIDS does not detect it (out of scope, with the reason) \\
\bottomrule
\end{tabularx}
\end{center}

\subsection{What TinyIDS does after a detection (the Responder)}

\begin{center}
\begin{tabularx}{\linewidth}{P{3.8cm}P{4.6cm}L}
\toprule
Action & When & What happens \\
\midrule
\textbf{Block} & sure attack (confidence $\geq$ 0.90 or a rule) from an outside machine & iptables drops
\emph{all} its packets for \textbf{30 s} \\
\textbf{Quarantine} & sure attack from one of \emph{our} NodeMCUs (it is compromised) & the same drop, but for
\textbf{120 s}; afterwards 30 s of grace so its reconnect does not trigger a new quarantine \\
\textbf{Rate-limit} & unsure attack (0.70--0.90, and the same sender flagged 2 windows in a row) & only packets above
\textbf{20 per second} are dropped, for 30 s; becomes a block if a sure alert follows \\
\textbf{Block MAC + fix ARP} & spoofing (the IP is fake, so blocking the IP is useless) & drop everything from the
attacker's network card (MAC) for 30 s; for ARP spoofing also write the correct MAC back into the Pi's ARP table \\
\bottomrule
\end{tabularx}
\end{center}

Every action ends by itself. The admin laptop's SSH session and open dashboard connection always stay open, so the
IDS can never lock out its own operator. The dashboard can switch blocking off (detect only) or release everything.

% ===============================================================================================================
\section{All attacks at a glance}

\begin{longtable}{P{4.3cm}P{3.5cm}P{3.1cm}P{3.6cm}}
\toprule
Attack & Detected by & Action & Status \\
\midrule
\endhead
\multicolumn{4}{l}{\textcolor{cFlood}{\textbf{Floods (DoS): one sender}}} \\
SYN flood & model + SYN cookies & block 30 s & live, 0.96 \\
UDP flood & model & block 30 s & live, 1.00 \\
ICMP (ping) flood & model & block 30 s & live, 1.00 \\
ACK flood (TCP flood) & model & block 30 s & live, 0.98 \\
PSH-ACK flood & model & block 30 s & live, 1.00 \\
RST-FIN flood & model & block 30 s & live, 1.00 \\
\multicolumn{4}{l}{\textcolor{cDDoS}{\textbf{Distributed floods (DDoS)}}} \\
Botnet DDoS (our NodeMCU bots) & model + FloodCounter rule & quarantine 120 s & live, TCP 0.96--0.97, UDP 1.00 \\
Mirai floods (3 kinds) & model (not trained on) & block / quarantine & dataset: 99.1--99.4\% \\
SynonymousIP flood & model on dataset rows; see card & see card & dataset: 99.9\% \\
\multicolumn{4}{l}{\textcolor{cSpoofFlood}{\textbf{Spoofed floods}}} \\
IP-spoofed flood (fake source IPs) & AddressWatcher rule & block MAC 30 s & Docker \\
\multicolumn{4}{l}{\textcolor{cFrag}{\textbf{Fragmentation floods}}} \\
ICMP fragmentation & model (not trained on) & block 30 s & dataset: 97.7\% \\
UDP fragmentation & model (not trained on) & block / rate-limit & dataset: 66.2\% \\
ACK fragmentation & model (not trained on) & block / rate-limit & dataset: 43.8\% \\
\multicolumn{4}{l}{\textcolor{cRecon}{\textbf{Reconnaissance (scans)}}} \\
Port scan & model + PortScanRule & block 30 s & live, 0.99 / 1.00 \\
Version / vulnerability scan & model & block 30 s & live (version scan), 0.99 \\
OS scan & model (named ``Flood'') & block 30 s & live, 0.96 \\
Ping sweep, host discovery & not detected (too few packets) & -- & not covered \\
\multicolumn{4}{l}{\textcolor{cBrute}{\textbf{Password attacks}}} \\
MQTT brute force & BruteForceRule & block 30 s & live, 1.00 \\
\multicolumn{4}{l}{\textcolor{cARP}{\textbf{Man in the middle}}} \\
ARP spoofing & AddressWatcher rule & block MAC 30 s + fix ARP & Docker \\
\multicolumn{4}{l}{\textcolor{cNone}{\textbf{Out of scope}}} \\
Web attacks, DNS spoofing, HTTP floods, SlowLoris & -- & -- & not covered \\
\bottomrule
\end{longtable}

% ===============================================================================================================
\section{Floods (DoS): one attacker}

\textbf{The idea of every flood:} send so many packets that the victim (or its network link) spends all its time on
garbage and has nothing left for real work. A \textbf{DoS} (Denial of Service) comes from one sender. The flood types
differ only in \emph{which} packets are sent, so they hurt different parts of the victim.

\textbf{How TinyIDS sees all of them:} a flood is many tiny, identical packets. In a 10-packet window that shows as
small and equal sizes (\file{Min}, \file{Max}, \file{Std} near 0) and one flag or protocol in 100\% of the packets.
The model calls it \textbf{Flood}; the SHAP reasons name the type; the FloodCounter rule counts the senders: one
sender $\rightarrow$ ``DoS flood''.

\begin{attack}{cFlood}{SYN flood}
\MODEL\ \KERNEL\ \LIVE
\fld{What it is}{A stream of TCP SYN packets: ``I want to open a connection'', thousands per second, often from fake
addresses.}
\fld{How it works}{For every SYN the victim reserves memory for a half-open connection and waits for the final ACK,
which never comes. The table of half-open connections fills up, and real clients can no longer connect.}
\fld{Why attackers do it}{Cheap and classic: one laptop with \texttt{hping3} is enough, and it attacks the place
every TCP service (MQTT, SSH, web) needs.}
\fld{Harm}{The broker refuses new connections: sensors that reconnect cannot publish any more; the admin cannot log in.}
\fld{How it is usually stopped}{SYN cookies (the kernel stops keeping state for half-open connections), rate limits
on SYNs, firewalls / DDoS protection in front of the server.}
\fld{How SYN cookies work}{Normally the server writes every SYN into a table and waits. With SYN cookies it
writes \textbf{nothing}: it hides the connection details (a short code made from the addresses, ports and a secret)
inside the sequence number of its SYN-ACK reply. A real client sends that number back in its final ACK, so the
server can rebuild the connection from it. A fake SYN never answers, so it costs the server no memory: the table
can no longer fill up.}
\fld{How TinyIDS stops it}{\textbf{Kernel:} SYN cookies are switched on at start, so the Pi survives even before the
alert. \textbf{Model:} Flood, live confidence \textbf{0.96}, SHAP reasons \file{Min}, \file{Std},
\file{ack\_flag\_number}. \textbf{Responder:} block the sender for 30 s.}
\end{attack}

\begin{attack}{cFlood}{UDP flood}
\MODEL\ \LIVE
\fld{What it is}{Large numbers of UDP packets to one or many ports.}
\fld{How it works}{UDP needs no connection, so the attacker just fires. The victim must look at every packet (and
for closed ports answer with an ICMP error); the Wi-Fi link fills up.}
\fld{Why attackers do it}{The easiest flood to send and to fake; high volume with no handshake.}
\fld{Harm}{Bandwidth and CPU are used up: real sensor messages are delayed or lost. A strong enough flood can even
push a small device off the Wi-Fi.}
\fld{How it is usually stopped}{Rate limits per source, dropping traffic to unused ports, upstream filtering.}
\fld{How TinyIDS stops it}{\textbf{Model:} Flood, live \textbf{1.00}, reasons \file{Protocol Type}, \file{Min},
\file{UDP}. \textbf{Responder:} block 30 s.}
\end{attack}

\begin{attack}{cFlood}{ICMP (ping) flood}
\MODEL\ \LIVE
\fld{What it is}{A storm of ICMP echo requests (``ping'').}
\fld{How it works}{The victim answers every ping with a reply, so it spends CPU and doubles the traffic on its own
link.}
\fld{Why attackers do it}{Works against almost any device (everything answers ping), needs no open port.}
\fld{Harm}{Slow or unreachable gateway; Wi-Fi airtime wasted.}
\fld{How it is usually stopped}{Limit or drop ICMP echo from outside, rate-limit ICMP.}
\fld{How TinyIDS stops it}{\textbf{Model:} Flood, live \textbf{1.00}, reasons \file{ICMP}, \file{Min},
\file{ack\_flag\_number}. \textbf{Responder:} block 30 s.}
\end{attack}

\begin{attack}{cFlood}{ACK flood (``TCP flood'' in the dataset)}
\MODEL\ \LIVE
\fld{What it is}{TCP packets with only the ACK flag, belonging to no real connection.}
\fld{How it works}{The victim (and every firewall on the way) must look each packet up in its connection table to
find out it belongs to nothing. ACK packets also slip past simple ``block new connections'' filters.}
\fld{Why attackers do it}{Gets past SYN-flood protections and stateless filters.}
\fld{Harm}{CPU spent on lookups, link filled.}
\fld{How it is usually stopped}{Stateful firewalls that drop ACKs without a connection, rate limits.}
\fld{How TinyIDS stops it}{\textbf{Model:} Flood, live \textbf{0.98}, reasons \file{Min}, \file{Max}, \file{Std}
(all packets the same tiny size). \textbf{Responder:} block 30 s.}
\end{attack}

\begin{attack}{cFlood}{PSH-ACK flood}
\MODEL\ \LIVE
\fld{What it is}{TCP packets with PSH and ACK set: ``here is data, deliver it at once''.}
\fld{How it works}{The PSH flag tells the receiver to pass the data straight to the application, so even more work
per packet than an ACK flood.}
\fld{Why attackers do it}{Looks like normal data traffic, so it is harder to filter by flags.}
\fld{Harm}{CPU and memory on the victim, link filled.}
\fld{How it is usually stopped}{Stateful filtering, rate limits.}
\fld{How TinyIDS stops it}{\textbf{Model:} Flood, live \textbf{1.00}, reasons \file{psh\_flag\_number}, \file{Min},
\file{Max}. \textbf{Responder:} block 30 s.}
\end{attack}

\begin{attack}{cFlood}{RST-FIN flood}
\MODEL\ \LIVE
\fld{What it is}{TCP packets with RST (reset) and FIN (finish) flags: ``close this connection''.}
\fld{How it works}{The victim must check every packet against its open connections to decide whether to tear one
down; a lucky guess can really reset a real connection.}
\fld{Why attackers do it}{Wastes the victim's lookups and can break real sessions.}
\fld{Harm}{CPU spent; real connections (a sensor's MQTT session) may be cut.}
\fld{How it is usually stopped}{Stateful firewalls, sequence-number checks, rate limits.}
\fld{How TinyIDS stops it}{\textbf{Model:} Flood, live \textbf{1.00}, reasons \file{fin\_flag\_number}, \file{Min},
\file{rst\_flag\_number}. \textbf{Responder:} block 30 s.}
\end{attack}

% ===============================================================================================================
\section{Distributed floods (DDoS)}

\textbf{DDoS} = the same floods, sent by many machines at once (a \emph{botnet}). Each one alone may look small; together
they are huge, and blocking one sender does not stop the attack.

\begin{attack}{cDDoS}{Botnet DDoS (our 3 NodeMCU bots)}
\MODEL\ \RULE{FloodCounter}\ \LIVE
\fld{What it is}{Hacked IoT devices (``bots'') that flood a target on command from the attacker's server. In our
testbed the 3 NodeMCUs carry bot firmware; the laptop's \file{bot.py} is the command server.}
\fld{How it works}{The attacker sends one command (``tcp 10''); every bot floods the target at the same time. Our bots
send UDP to port 5683 or open TCP connections to port 1884 for at most 30 s.}
\fld{Why attackers do it}{IoT devices are cheap to hack (default passwords, no updates) and there are millions of
them. Motives: extortion (``pay or we keep you offline''), revenge, hiding another attack, renting the botnet out.}
\fld{Harm}{Much more traffic than one sender; and the attackers are \emph{our own} devices, which are compromised.}
\fld{How it is usually stopped}{Upstream DDoS scrubbing, rate limits, and finding and cleaning the infected devices.}
\fld{How TinyIDS stops it}{\textbf{Model:} Flood for each bot (live: TCP \textbf{0.96--0.97}, UDP \textbf{1.00}),
reasons \file{syn\_flag\_number}, \file{fin\_flag\_number}, \file{rst\_flag\_number}. \textbf{FloodCounter rule:}
2 or more senders flooding within 10 s $\rightarrow$ ``DDoS flood (3 senders)''. \textbf{Responder:} the senders are
our own sensors, so they are \textbf{quarantined for 120 s} (live: all 3 quarantined).}
\end{attack}

\begin{attack}{cDDoS}{Mirai floods (greeth, greip, udpplain)}
\MODEL\ \DATA
\fld{What it is}{Floods sent by \textbf{Mirai}, the famous IoT botnet (2016) that infected cameras and routers through
default passwords and took large websites offline.}
\fld{How it works}{\file{greeth} and \file{greip} send GRE-tunnelled packets (Ethernet or IP inside GRE);
\file{udpplain} is a plain, fast UDP flood.}
\fld{Why attackers do it}{Ready-made, open-source botnet code; huge numbers of weak IoT devices.}
\fld{Harm}{Same as a DDoS; plus every infected device is under the attacker's control.}
\fld{How it is usually stopped}{Change default passwords, update firmware, block telnet from outside, DDoS protection.}
\fld{How TinyIDS stops it}{Not in training (we cannot run real Mirai), but the floods look like our floods: on the
dataset's test rows the model calls them attacks \textbf{99.1--99.4\%} of the time $\rightarrow$ block (or
quarantine if it is one of our sensors). Our bot firmware imitates Mirai's command-and-flood behaviour safely.}
\end{attack}

\begin{attack}{cDDoS}{SynonymousIP flood}
\MODEL\ \DATA\ \NOTEST
\fld{What it is}{A flood label in CICIoT2023 (\file{DDoS-SynonymousIP\_Flood}). It is commonly described as a TCP SYN
flood whose \emph{fake source address is the victim's own IP address} (like the old ``LAND'' attack). We did not
check this against the dataset's original packet captures.}
\fld{How it works}{The victim receives connection requests that seem to come from itself, so it answers itself and can
get stuck in a loop; with many senders it is also a plain SYN flood.}
\fld{Why attackers do it}{Confuses devices with weak network stacks; the fake source hides the real sender.}
\fld{Harm}{Like a SYN flood, plus possible loops or crashes in poorly written IoT firmware.}
\fld{How it is usually stopped}{Drop packets that claim to come from your own address (Linux does this by default:
``martian'' source), SYN cookies, anti-spoofing filters at the router.}
\fld{How TinyIDS stops it}{On the dataset's rows the model calls it an attack \textbf{99.9\%} of the time. Live, the
detector ignores packets that claim to come from the Pi itself, so the model would not see this exact variant; the
Linux kernel's own check and the SYN cookies are the protection. If the fake sources are many \emph{different}
addresses, the AddressWatcher rule (next section) catches it.}
\end{attack}

% ===============================================================================================================
\section{Spoofed floods}

\begin{attack}{cSpoofFlood}{IP-spoofed flood (random fake source addresses)}
\RULE{AddressWatcher}\ \DOCKER
\fld{What it is}{Any flood (SYN, UDP, ICMP\ldots) where every packet carries a different, made-up source IP
(\texttt{hping3 --rand-source}).}
\fld{How it works}{The victim cannot tell who is really sending, and blocking ``the attacker's IP'' is useless: the next
packet has a new one.}
\fld{Why attackers do it}{Hide the real sender, defeat per-IP blocking, and make replies hit innocent addresses.}
\fld{Harm}{Same as the flood type, but harder to stop.}
\fld{How it is usually stopped}{Anti-spoofing filters at routers (BCP 38 / reverse-path checks), DDoS scrubbing.}
\fld{How TinyIDS stops it}{The model cannot help (no sender ever fills a 10-packet window, every packet has a new
address). \textbf{AddressWatcher rule:} one network card (MAC) using \textbf{more than 20 source IPs within 5 s}
$\rightarrow$ IP spoofing. \textbf{Responder:} the IP is fake, so the \textbf{MAC is blocked} for 30 s. Works inside
the local network (MAC addresses do not cross routers), which is our case.}
\end{attack}

% ===============================================================================================================
\section{Fragmentation floods}

\textbf{Fragmentation:} a packet too big for the network is cut into pieces (fragments) and the receiver puts them
back together. To do that it must keep the pieces in memory until all have arrived. Fragmentation floods abuse this:
they send pieces that are expensive to reassemble or never complete.

\textbf{For all three:} not in our training and never tested on our testbed (\texttt{hping3 -f} could send them).
Common protection: limits on the reassembly memory and time (Linux has them), dropping fragments where they are not
needed, rate limits.

\begin{attack}{cFrag}{ICMP fragmentation flood}
\MODEL\ \DATA
\fld{What it is}{Large ping packets cut into fragments, sent in huge numbers.}
\fld{How it works}{The victim must buffer and reassemble every fragment set before it can even answer the ping.}
\fld{Why attackers do it}{More work for the victim per packet than a plain ping flood.}
\fld{Harm}{Memory for reassembly fills up, CPU wasted, link filled.}
\fld{How TinyIDS stops it}{It looks very much like an ICMP flood: the model calls the dataset's rows an attack
\textbf{97.7\%} of the time $\rightarrow$ block 30 s.}
\end{attack}

\begin{attack}{cFrag}{UDP fragmentation flood}
\MODEL\ \DATA
\fld{What it is}{Large UDP packets sent as fragments.}
\fld{How it works}{Same reassembly work; the fragments carry no UDP header after the first piece, so port-based
filters cannot judge them.}
\fld{Why attackers do it}{Bypasses port filters; consumes reassembly memory.}
\fld{Harm}{Memory and bandwidth.}
\fld{How TinyIDS stops it}{Partly: the model flags the dataset's rows \textbf{66.2\%} of the time (block, or
rate-limit when unsure). Its fragments look less like our UDP floods than the ICMP ones do.}
\end{attack}

\begin{attack}{cFrag}{ACK fragmentation flood}
\MODEL\ \DATA
\fld{What it is}{TCP ACK packets cut into fragments.}
\fld{How it works}{Combines the ACK flood (connection-table lookups) with reassembly work; the pieces are big, so few
packets fill the link.}
\fld{Why attackers do it}{High bandwidth with few packets; passes simple flag filters.}
\fld{Harm}{Bandwidth and memory.}
\fld{How TinyIDS stops it}{Weakly: the model flags the dataset's rows only \textbf{43.8\%} of the time. A known gap: to
cover it, record it with \texttt{hping3 -A -f} and add it to training.}
\end{attack}

% ===============================================================================================================
\section{Reconnaissance (scans)}

\textbf{Reconnaissance} is the attacker looking around before the real attack: which devices exist, which ports are
open, which software and versions run. It does no damage by itself, but it is usually the \emph{first step}, so
catching it early stops the attack before it starts.

\begin{attack}{cRecon}{Port scan}
\MODEL\ \RULE{PortScanRule}\ \LIVE
\fld{What it is}{Trying many ports of the target to see which are open (\texttt{nmap}).}
\fld{How it works}{A SYN is sent to port after port; an answer means ``open''. Fast and quiet: no full connection.}
\fld{Why attackers do it}{To find a door: an open MQTT broker, SSH, a forgotten web panel.}
\fld{Harm}{None directly; it reveals what to attack next.}
\fld{How it is usually stopped}{Firewall all unused ports, scan detection in IDSs, rate limits.}
\fld{How TinyIDS stops it}{\textbf{Model:} Recon, live \textbf{0.99}, reasons \file{AVG}, \file{Tot size},
\file{Max}. \textbf{PortScanRule:} one sender touching \textbf{more than 20 ports within 5 s} $\rightarrow$ Recon
(live \textbf{1.00}). Both fired. \textbf{Responder:} block 30 s, so the scan sees nothing more.}
\end{attack}

\begin{attack}{cRecon}{Version / vulnerability scan}
\MODEL\ \LIVE
\fld{What it is}{After finding open ports, asking each service what software and version it runs
(\texttt{nmap -sV}), or testing known weaknesses.}
\fld{How it works}{Sends probes that make services reveal their name and version, then compares with lists of known
vulnerabilities.}
\fld{Why attackers do it}{To pick an exploit that works on exactly this version.}
\fld{Harm}{Reveals weaknesses; some probes can crash fragile devices.}
\fld{How it is usually stopped}{Hide version banners, keep software updated, IDS scan detection.}
\fld{How TinyIDS stops it}{\textbf{Model:} Recon, live \textbf{0.99} on our version scan (reasons \file{AVG},
\file{Tot size}, \file{Max}) $\rightarrow$ block 30 s. On the dataset's vulnerability-scan rows the model is weaker
(36\% flagged).}
\end{attack}

\begin{attack}{cRecon}{OS scan}
\MODEL\ \LIVE
\fld{What it is}{Guessing the target's operating system from small details of how it answers odd packets
(\texttt{nmap -O}).}
\fld{How it works}{Unusual flag combinations and probes; every OS answers slightly differently.}
\fld{Why attackers do it}{Knowing ``Linux on a Raspberry Pi'' tells them which exploits to try.}
\fld{Harm}{Reveals information; no damage by itself.}
\fld{How it is usually stopped}{Firewalls that drop odd probes, IDS scan detection.}
\fld{How TinyIDS stops it}{Live: the model called it \textbf{Flood} (0.96, reasons \file{Min}, \file{Std},
\file{syn\_flag\_number}): the wrong name, because its probes look like a small SYN flood, but it was
\textbf{blocked at once}. On the dataset's OS-scan rows the model is weak (25\% flagged).}
\end{attack}

\begin{attack}{cNone}{Ping sweep and host discovery}
\NOTCOV
\fld{What it is}{Finding which devices are on the network by pinging every address.}
\fld{Why it is not covered}{Each device receives only one or two packets, so no sender ever sends the Pi the 10
packets a window needs: the model never sees it. It only learns that the Pi exists, which is not secret on a shared
Wi-Fi.}
\end{attack}

% ===============================================================================================================
\section{Password attacks}

\begin{attack}{cBrute}{MQTT brute force}
\RULE{BruteForceRule}\ \LIVE
\fld{What it is}{Trying password after password on the MQTT broker (\texttt{mosquitto\_pub} with wrong passwords).}
\fld{How it works}{Each try is a CONNECT; the broker answers with a CONNACK that says ``refused''. The attacker keeps
going until one is accepted.}
\fld{Why attackers do it}{With the broker password they can read every sensor value and \emph{publish fake ones}, or
send commands to devices.}
\fld{Harm}{Full control of the IoT messages: fake readings, hidden real alarms, commands to actuators.}
\fld{How it is usually stopped}{Strong passwords, lockout after failed logins, TLS client certificates.}
\fld{How TinyIDS stops it}{The model cannot see it (a login is a normal-looking packet). \textbf{BruteForceRule:}
the detector reads the broker's answers; \textbf{5 refused logins within 10 s} from one client $\rightarrow$
BruteForce (live \textbf{1.00}). \textbf{Responder:} block 30 s.}
\end{attack}

% ===============================================================================================================
\section{Man in the middle}

\begin{attack}{cARP}{ARP spoofing}
\RULE{AddressWatcher}\ \DOCKER
\fld{What it is}{The attacker tells the network ``the Pi's address belongs to \emph{my} network card''
(\texttt{arpspoof}), so traffic meant for one device goes to the attacker.}
\fld{How it works}{Devices find each other's network card with ARP and trust any answer. Fake ARP answers poison their
tables.}
\fld{Why attackers do it}{To sit in the middle: read every sensor message, change them in flight, or drop them.}
\fld{Harm}{Spying, fake data, cut-off devices; the victims notice nothing.}
\fld{How it is usually stopped}{Static ARP entries for important devices, ``dynamic ARP inspection'' on managed
switches, encryption (TLS) so stolen traffic is useless.}
\fld{How TinyIDS stops it}{\textbf{AddressWatcher rule:} it remembers which MAC owns which IP; an ARP message that
gives a known IP a \emph{different} MAC $\rightarrow$ ARP spoofing. \textbf{Responder:} block the attacker's MAC for
30 s and \textbf{write the correct MAC back} into the Pi's ARP table. Tested in Docker; not yet on the real Pi.}
\end{attack}

% ===============================================================================================================
\section{Not covered, and why}

\begin{attack}{cNone}{Web attacks (SQL injection, XSS, command injection, backdoor, browser hijacking, uploading)}
\NOTCOV
\fld{What they are}{Attacks on web applications: bad input that makes a website run the attacker's code or queries.}
\fld{Why not covered}{The Pi runs no web server, so there is nothing to attack and nothing to test. Their packets also
look like normal traffic in our 14 features (model: 0--5\% on the dataset's rows).}
\end{attack}

\begin{attack}{cNone}{DNS spoofing}
\NOTCOV
\fld{What it is}{Fake answers to name lookups, sending a device to the attacker's server.}
\fld{Why not covered}{Our devices use fixed IP addresses (no lookups); the model flags the dataset's rows only 15\%.}
\end{attack}

\begin{attack}{cNone}{HTTP floods and SlowLoris}
\NOTCOV
\fld{What they are}{Floods of web requests, or (SlowLoris) many web connections kept open by sending a request very
slowly.}
\fld{Why not covered}{They need a web server to attack; the Pi has none. (The model still flags 40--61\% of the
dataset's HTTP-flood rows.)}
\end{attack}

% ===============================================================================================================
\section{The layers of defence, in one table}

\begin{center}
\begin{tabularx}{\linewidth}{P{3.6cm}P{4.6cm}L}
\toprule
Layer & What it does & Against \\
\midrule
Linux kernel & SYN cookies; drops packets claiming the Pi's own address & SYN floods, SynonymousIP-style floods \\
Model (LightGBM) & judges the shape of every 10-packet window & all floods, scans, DDoS (per sender) \\
FloodCounter & counts flooding senders & DoS vs DDoS \\
PortScanRule & counts ports per sender & fast port scans \\
BruteForceRule & counts refused MQTT logins & password guessing \\
AddressWatcher & checks IP $\leftrightarrow$ MAC & ARP spoofing, spoofed-source floods \\
Responder (iptables) & block / quarantine / rate-limit / MAC block, all time-limited & everything detected above \\
Dashboard & shows every alert with its reasons; IPS on/off, release all & the human in the loop \\
\bottomrule
\end{tabularx}
\end{center}

\end{document}
